%% kapa: exam=1 round=31 year=2020 period=1 form=A subject=real_estate subject_ko=부동산학원론 seq=37 number=117
%% source: Q-Net 감정평가사 기출문제 (cst00309) artlSeq=5207219 file=extracted/1008/5207219/A형/A형/2020년도 제31회 감정평가사 1차 시험_1교시_A형.pdf
%% conversion: kordoc 4.24.51 -> tools/gen1.py (auto); review=proofread-sonnet; flags=none
\begin{question}{117}
공인중개사법령에 관한 설명으로 옳은 것은?
\begin{choices}
  \item 공인중개사법에 의한 공인중개사자격을 취득한 자를 개업공인중개사라고 말한다.
  \item 선박법 및 선박등기법에 따라 등기된 20톤 이상의 선박은 공인중개사법에 의한 중개 대상물이다.
  \item 개업공인중개사에 소속된 공인중개사인 자로서 중개업무를 수행하는 자는 소속공인중개사가 아니다.
  \item 중개업은 다른 사람의 의뢰에 의하여 일정한 보수를 받고 중개를 업으로 행하는 것을 말한다.
  \item 중개보조원이란 공인중개사가 아닌 자로서 중개업을 하는 자를 말한다.
\end{choices}
\end{question}
\answer{4}
